\documentclass[12pt]{article}
\usepackage[margin=1in, paperwidth=8.5in, paperheight=11in]{geometry}
\usepackage{ifpdf,amsmath, amssymb, comment, color, graphicx, stmaryrd,setspace,enumerate,tikz, fancyhdr, wrapfig, textcomp, mathptmx, siunitx}
\usepackage{pdfpages}
\usepackage{enumerate}

\usepackage{hyperref}
\hypersetup{
    colorlinks=true,
    urlcolor=blue,
}

\setlength{\headheight}{14.5pt}
\newcommand{\Q}{\mathbb{Q}}
\newcommand{\R}{\mathbb{R}}
\newcommand{\Z}{\mathbb{Z}}

\newenvironment{red}{\color{red}}{\ignorespacesafterend}


\newcommand{\blue}[1]{\textcolor{blue}{#1}}
\newcommand{\green}[1]{\textcolor{green}{#1}}
\renewcommand{\section}[1]{\begin{center} \textbf{#1} \\\end{center}}
%
\setlength{\parindent}{0in}
%\oddsidemargin=-.25in
\allowdisplaybreaks
\pagestyle{fancy}
\renewcommand{\headrulewidth}{0pt}
\lhead{MATH 202}
\rhead{Fall 2024}
%\lfoot{\copyright\ CLEAR Calculus 2010}
\cfoot{}
\renewcommand{\thefootnote}{*} 
\hyphenpenalty=10000
\begin{document}
%


%\onehalfspacing
\allowdisplaybreaks
%##################################################################
\section{Physics applications of definite integrals}
Here are a bunch of situations where we want to calculate something by multiplying some stuff together, but it doesn't work because one of those things isn't constant -- it's some function of something else -- so we have to slice-approximate-integrate.

\hrulefill

\begin{enumerate}
\item The kinetic energy of an object with mass \(m\) and constant speed \(v\) is \(K = \frac{1}{2} mv^2\) (measured in Joules, \(\unit{\J} = \unit{\kg \cdot \m^2 / \s ^2} \)). Picture the second hand of a clock: it's a rod that is 0.1 m long with a total mass of 0.03 kg, and it is rotating around one of its ends at a rate of one revolution per minute. I'd like to know the total kinetic energy of this object.
\begin{enumerate}[(a)]
\item Draw a picture.
\item Which quantity in the formula for kinetic energy isn't constant throughout your picture? What quantity does it depend on? Label that quantity in your picture.

\begin{red}
Velocity is a nonconstant function of radius, $v(r)$. 
\end{red}
\item Draw a representative small slice where that quantity is approximately constant. Label that slice with a \(\Delta\)\rule{0.5cm}{0.15mm}.
\item Write an equation for \(\Delta K\) for that slice.

\begin{red}
    Note that the uniform density of the rod is $\frac{\SI{0.03}{kg}}{\SI{0.1}{m}} = \SI{0.3}{kg/m}$, which we'll call $\rho$.
    \[\Delta K = 1/2m(v(r))^2 = 1/2\cdot (\rho\cdot\Delta r)\cdot\left(\frac{2\pi}{60}\cdot r\right)^2, \]
where $r$ is the distance from the end of the rod to the representative slice.
\end{red}
\item Write a definite integral that gives the total kinetic energy \(K\) of the entire second hand.

\begin{red}
    \[K_{total} = \int_0^{0.1} \frac{1}{2}\cdot 0.3 \left(\frac{2\pi r}{60}\right)^2\, dr\]
\end{red}
\item Evaluate the integral, assign it the proper units, and explain the meaning of your result.

\begin{red}
    \begin{align*} 
        \int_{0}^{0.1} \frac{1}{2} \cdot 0.3 \left(\frac{2\pi r}{60} \right)^2\ dr 
        & = 0.001645 \int_0^{0.1} r^2\, dr \\ 
        & = 0.000548 r^3 \bigg|_0^{0.1}\\ 
        & = \SI{5.483e-7}{J}
    \end{align*}
    So, the rod has a total energy of $\SI{5.483e-7}{J}$.
\end{red}
\item Describe the meaning of each factor of your integral, and give its units.

\begin{red}
    \[\int_0^{0.1} \underbrace{\frac{1}{2} \left(\underbrace{\frac{2\pi r}{60}}_{\text{Velocity of a slice at radius \(r\) (m/s)}\quad }\right)^2 \underbrace{\overbrace{0.3}^{\text{density (kg/m)}}\cdot \overbrace{dr}^{\text{thickness of slice (m)}}}_{\quad \text{mass of a slice of}\ ``\text{thickness \(dr\)" (kg)} } }_{ \text{kinetic energy of a slice of}\ ``\text{thickness \(dr\)" (J)} } \]
\end{red}
\end{enumerate}

\item Picture an oil slick on the surface of the ocean. It's probably roughly a circle, and the oil is probably denser in the center than it is at the edges. I'd like to know the total mass of oil in this oil slick; let's say that the density of oil at a distance \(r\) meters from the center of the slick is given by \[\rho(r) = \dfrac{50}{1 + r}\, \unit{\kg / \m^2}.\]
\begin{enumerate}[(a)]
    \item Draw a picture.
    \item Which quantity in the formula for mass isn't constant throughout your picture? What quantity does it depend on? Label that quantity in your picture.

    \begin{red}
        The issue is that density is a nonconstant function of radius, $\rho(r)$.
    \end{red}
    \item Draw a representative small slice where that quantity is approximately constant. Label that slice with a \(\Delta\)\rule{0.5cm}{0.15mm}.
    \item Write an equation for \(\Delta M\) for that slice.

    \begin{red}
        \[\Delta M = \rho(r)\cdot \Delta A = \frac{50}{1+r}\cdot (2\pi r\Delta r) \]
    \end{red}
    \item Write a definite integral that gives the total mass \(M\) of the entire oil slick; let's say that the oil slick is 20 km across.

    \begin{red}
        If the oil slick is 20 km across, then the radius is 10 km = 10,000 m.
        \[M_{total} = \int_0^{10000} \frac{50}{1+r}\cdot 2\pi r\, dr\]
    \end{red}
    \item Evaluate the integral, assign it the proper units, and explain the meaning of your result. Hint: \(\frac{r}{1+r} = \frac{r+1-1}{1+r} = 1 - \frac{1}{1+r}\).
    \begin{red}
        \begin{align*} 
            \int_0^{10000} \frac{50}{1 + r} 2\pi r\, dr 
            & = 100\pi \int_0^{10000} \frac{r}{1 + r}\ dr \\ 
            &  = 100\pi (r - \ln(1 + r)) \bigg|_0^{10000}\\ 
            & = 100\pi(10000 - \ln(10001))\\ & = \SI{3.1386e6}{kg}
        \end{align*}
        So, the total mass of the slick is $\SI{3.1386e6}{kg}.$
    \end{red}
    \item Describe the meaning of each factor of your integral, and give its units.
    \begin{red}
        \[\int_0^{10000} \underbrace{ \underbrace{\left(\frac{50}{1 + r}\right)}_{\text{Density at radius \(r\) (kg/m\(^2\))}\quad } \underbrace{2\pi r\, dr}_{\quad \text{Area of a ring at radius \(r\) with}\ ``\text{thickness \(dr\)" (m\(^2\))} } }_{ \text{Mass of oil in a ring at radius \(r\) with}\ ``\text{thickness \(dr\)" (kg)} }\]
    \end{red}
\end{enumerate}

\item The force of gravity that the earth exerts on an object diminishes as the object gets further away from the earth. Specifically, according to Newton, the force of gravity is governed by an ``inverse square law:'' for some constant \(k\) that depends on the masses of the earth and the object, the gravitational force at distance \(r\) from the center of the earth is given by 
\[ F(r) = \frac{k}{r^2}\, \unit{\N}.\]
The energy required to move an object a distance \(d\) is \(E = F\cdot d\), measured in Joules, \(\unit{\J} = \unit{\N \cdot \m}\), if the force is constant over the distance \(d\).  I'd like to know the total energy necessary to launch a \href{https://en.wikipedia.org/wiki/CubeSat}{CubeSat} with mass 2 kg all the way to the moon.
\begin{enumerate}[(a)]
    \item Draw a picture.
    \item Which quantity in the formula for energy isn't constant throughout your picture? What quantity does it depend on? Label that quantity in your picture.

    \begin{red}
        The issue is that the gravitational force of the earth on the satellite is a nonconstant function of $r$, the distance from the center of the earth to the satellite: $F(r)$.
    \end{red}
    \item Draw a representative small slice where that quantity is approximately constant. Label that slice with a \(\Delta\)\rule{0.5cm}{0.15mm}.
    \item Write an equation for \(\Delta E\) for that slice.
    \begin{red}
        \[\Delta E = F(r)\cdot d = \frac{k}{r^2}\cdot\Delta r\]
    \end{red}
    \item Write a definite integral that gives the total energy \(E\) needed for the entire launch. The distance from the surface of the earth to the surface of the moon is about \(362{,}000\unit{\km}\).
    \begin{red}
        \[E_{total} = \int_{6371000}^{368371000}\frac{k}{r^2}\,dr\]
        \textbf{Careful:} notice that the variable of integration $r$ is the distance from the \textbf{center} of the earth, not from the \textbf{surface} of the earth.
    \end{red}
    \item Let's figure out the right value for \(k\). Our CubeSat has mass 2 kg, and at the surface of the earth (\(r = \qty{6371}{\km}\)), the force of gravity on the CubeSat is \(\qty{2}{\kg}\cdot\qty{9.8}{\m/\s^2} = \qty{19.6}{\N}\).
    \begin{red}
        \[19.6 = F(6371000) = \frac{k}{(6371000)^2} \text{, so } k = (6371000)^2\cdot 19.6 \approx \num{7.956e14}. \]
    \end{red}
    \item Evaluate the integral, assign it the proper units, and explain the meaning of your result.

    \begin{red}
        I'm going to leave $k$ as a letter until the very end.
        \begin{align*}
            E_{total} &= \int_{6371000}^{368371000}\frac{k}{r^2}\,dr = \frac{-k}{r}\bigg|_{6371000}^{368371000} 
            = -k\left(\frac{1}{368371000} - \frac{1}{6371000}\right) \\
            &= -k\cdot\num{-1.54247e-7} = \num{7.956e14}\cdot\num{1.54203e-7} = \SI{1.227e8}{J}.
            \end{align*}
            So, the total energy required to move the CubeSat to the moon is  $\SI{1.227e8}{J}$.
    \end{red}
    \item Describe the meaning of each factor of your integral, and give its units.
    \begin{red}
        \[\int_{6371000}^{368371000} \underbrace{ \underbrace{\frac{k}{r^2} }_{\text{Force required at \(r\) (N)}\quad } \underbrace{dr}_{\quad \text{Distance moved (m)} } }_{ \text{Energy to move the 1 kg payload a distance \(dr\) (J)}} \]
    \end{red}
\end{enumerate}

\item The energy required to move an object a distance \(x\) while exerting a constant force \(F\) is \(E = F\cdot x\).  Suppose that you have two magnets and a wire. One magnet is attached to the end of the wire and the other can slide along the wire. If the magnets are arranged so that they repel each other, then it will require force to push the movable magnet toward the fixed magnet. The amount of force needed to move the magnet increases as the two get closer together. In fact, the force at a distance \(x\) is proportional to \(1/x^2\); that is, for some constant $k$, the force at a distance $x$ is $k\cdot 1/x^2$. I'd like to know the total energy needed to push the movable magnet a certain distance toward the fixed magnet.
\begin{enumerate}[(a)]
    \item Draw a picture.
    \item Which quantity in the formula for energy isn't constant throughout your picture? What quantity does it depend on? Label that quantity in your picture.

    \begin{red}
        The issue is that force is a nonconstant function of distance, $F(x)$.
    \end{red}
    \item Draw a representative small slice where that quantity is approximately constant. Label that slice with a \(\Delta\)\rule{0.5cm}{0.15mm}.
    \item Write an equation for \(\Delta E\) for that slice.
    \begin{red}
        \[\Delta E = F(x)\cdot d =  \frac{k}{x^2}\cdot\Delta x\]
    \end{red}
    \item Write a definite integral that gives the total energy \(E\) needed to push the magnet from 5 cm away to 3 cm away. 
    \begin{red}
        \[E_{total} = \int_{0.03}^{0.05} \frac{k}{x^2}\,dx\]
    \end{red}
    \item Evaluate the integral, assign it the proper units, and explain the meaning of your result. Your answer will involve the constant \(k\), which we don't have enough information to figure out. (What information would we need?)
    \begin{red}
        \[E_{total} = \int_{0.03}^{0.05} \frac{k}{x^2}\,dx = \frac{-k}{x}\bigg|_{0.03}^{0.05} = k\cdot\SI{13.33}{J}\]
        So, it takes 13.33$k$ J (not to be confused with 13.33 kilojoules!) to move the magnet this distance.
    \end{red}
    \item Describe the meaning of each factor of your integral, and give its units.
    \begin{red}
        \[\int_{0.03}^{0.05} \underbrace{ \underbrace{\frac{k}{x^2}}_{\text{Force required to move magnet at distance \(x\) (N)}\quad } \underbrace{dx}_{\text{Distance moved (m)} } }_{ \text{Energy required to move the magnet a distance \(dx\) when the magnets are \(x\) m apart (J)} }\]
    \end{red}
\end{enumerate}

\item Consider the atmosphere: it's quite thick at ground level, but thins out as you get higher and higher, until you get to space and it's basically nothing. I'd like to know the total mass of the atmosphere; it turns out that a good model for the density of the atmosphere as a function of height \(h\) (measured in meters up from sea level) is
\[ \rho(h) = 1.28 e^{-0.000124 h}\, \unit{\kg/\m^3}.\]
\begin{enumerate}[(a)]
    \item Draw a picture.
    \item Which quantity in the formula for mass isn't constant throughout your picture? What quantity does it depend on? Label that quantity in your picture.
    
    \begin{red}
        The issue is that density is a nonconstant function of height, $\rho(h)$.
    \end{red}
    \item Draw a representative small slice where that quantity is approximately constant. Label that slice with a \(\Delta\)\rule{0.5cm}{0.15mm}.
    \item Write an equation for \(\Delta M\) for that slice. \textbf{Be careful:} your equation will probably depend on \(r\), which is not the same as \(h\).

    \begin{red}
        First, since the slice is a very thin spherical "shell", we can approximate its volume as its surface area times its thickness: $\Delta V \approx 4\pi r^2\cdot \Delta r$. (Think about painting a basketball.)
        \[\Delta M = \rho(h) \cdot \Delta V = 1.28 e^{-0.000124 h} \cdot 4\pi r^2 \Delta r.\]
        Now, I'm irritated that this involves both $h$ and $r$, so note that $h = r - 6300000$ (since the radius of the earth is about $\SI{6300}{km}$):
        \[\Delta M = 1.28 e^{-0.000124 (r - 6300000)} \cdot 4\pi r^2\Delta r.\]
        (An alternative approach is to replace $r$ by $h + 6300000$: \\ 
        $\Delta M = 1.28 e^{-0.000124 h} \cdot 4\pi (h+6300000)^2 \Delta h.$)
    \end{red}
    \item Write a definite integral that gives the total mass \(M\) of the atmosphere all the way out to the ``Kármán line'' of \(\qty{83.8}{\km}\) above sea level, which is the highest height theoretically achievable by an aircraft. (Note: the earth's radius is about \(\qty{6300}{km}\).)

    \begin{red}
        Note that $h = 0$ corresponds to $r = 6300000$ and $h = 83800$ corresponds to $r = 6383800$.
    \[M_{total} = \int_{r=6300000}^{6383800} 1.28 e^{-0.000124 (r - 6300000)} \cdot 4\pi r^2\,dr\] or, equivalently,
    \[M_{total} = \int_{h=0}^{83800} 1.28 e^{-0.000124 h} \cdot 4\pi (h+6300000)^2\, dh \]
    \end{red}
    \item Evaluate the integral (use Wolfram Alpha or something), assign it the proper units, and explain the meaning of your result.

    \begin{red}
        Either integral gives the same value: the mass of the atmosphere out to the Kármán line is $\qty{5.162e18}{kg}$.
    \end{red}
    \item Describe the meaning of each factor of your integral, and give its units.
    \begin{red}
        \[\int_{r=6300000}^{6383800} 
    \underbrace{ 
        \underbrace{ 1.28 e^{-0.000124(r-6300000)} }_
            { \stackrel{ \text{Density of the earth's atmosphere}}
                       { \text{at altitude $r$ (kg/m\(^2\))   }\quad } } 
        \underbrace{ 4\pi r^2\ dr }_
            { \stackrel{ \text{Volume of a thin shell of the earth's} }
                       { \text{atmosphere at altitude $r$ (m\(^3\))} } } }
    _{ \text{Mass of a spherical shell with thickness \(dr\) of the earth's atmosphere 
            at altitude $r$ (kg)} }\]
    \end{red}
\end{enumerate}

\item The gravitational attraction between two \textbf{particles} of mass \(m_1\) and \(m_2\) at a distance \(r\) apart is
\[F(r) = \frac{Gm_1 m_2}{r^2}.\]
Picture a thin uniform rod with mass \(M\) and length \(\ell\), and then a particle of mass \(m\) that's \(a\) meters away from the end of the rod. (I am imagining a pool cue hitting a pool ball.) I'd like to know the total gravitational force between the rod and the particle; \textbf{the problem is that the rod is not a particle.}
\begin{enumerate}[(a)]
    \item Draw a picture.
    \item Which quantity in the formula for force isn't constant throughout your picture? What quantity does it depend on? Label that quantity in your picture.

    \begin{red}
        Since the rod is not a particle, the issue is that there isn't a consistent distance between a piece of the rod and the ball. (Where do you measure from?) Therefore we must cut up the rod into a bunch of little pieces (the fancy word is ``discretize'' the rod) of mass $\Delta M$, and then $r$ is nonconstant.
    \end{red}
    \item Draw a representative small slice where that quantity is approximately constant. Label that slice with a \(\Delta\)\rule{0.5cm}{0.15mm}.
    \item Write an equation for \(\Delta F\) for that slice.

    \begin{red}
        We'll cut the rod into pieces of length $\Delta r$. Since the rod has uniform density (which is given by dividing the total mass of the rod by the total length of the rod, $\rho = \dfrac{M}{\ell}\si{kg/m}$), the mass of a small slice $\Delta M = \rho \cdot \Delta r = \dfrac{M\cdot \Delta r}{\ell}$. Therefore,
        \[\Delta F = \dfrac{G\cdot \Delta M \cdot m}{r^2} = \dfrac{1}{r^2} \cdot \left( G \cdot \left(\dfrac{M\cdot \Delta r}{\ell}\right)\cdot m \right)  = \dfrac{G\cdot M \cdot m}{\ell\cdot r^2} \Delta r.\]
    \end{red}

    \item Write a definite integral that gives the total force \(F\) between the rod and the particle.

    \begin{red}
        The first slice of the rod is located at a distance $a$ away from the ball; the last slice of the rod is located at a distance $a+\ell$ away from the ball.
        \[F_{total} = \int_{r=a}^{a+\ell}\, \dfrac{G\cdot M \cdot m}{\ell\cdot r^2}\, dr = \dfrac{GMm}{\ell} \int_{r=a}^{a+\ell} \dfrac1{r^2}\,dr.\]
    \end{red}
    \item Evaluate the integral (your answer will involve lots of letters), assign it the proper units, and explain the meaning of your result.
    \begin{red}
        \begin{align*}
            F_{total} &= \dfrac{GMm}{\ell} \int_{r=a}^{a+\ell} \dfrac1{r^2}\,dr \\
            &= \frac{GMm}{\ell}\cdot\left[\frac{-1}{r}\right|_a^{a+\ell} 
             = \frac{GMm}{\ell}\left(\frac{1}{a} - \frac{1}{a+\ell}\right) \\
            &= \frac{GMm}{\ell}\cdot\left(\frac{a+\ell-a}{a(a+\ell)}\right) 
             = \frac{GMm}{a(a+\ell)} \si{N}.
        \end{align*}
    \end{red}
\end{enumerate}

\end{enumerate}
%#####################################d############################
\end{document}